\def\ChapterCode{ARZ}
\chapter
[Arzneimittel]
{Arzneimittel}
\label{chp:ARZ}




\ChapterIntro
{%

}
{%
\begin{description}
  \item[Maintainer:] Sebastian Gabriel
  \item[Autoren:] Diverse
  \item[Reviewer:] Standard-Reviewprozess
  \item[Version:] {\VarDocVersionTypeName} 
  ({\VarDocVersionTypeDescription})
  \item[SHA1:] {\DocGitCommit}
\end{description}

{\TxTeilkapitelAASS}
}

\opt{STATUS-EXPERIMENTAL}{2\Mg

2\Mg*

2\,\Mg

2\Ml

2\Ml*}


\section{Umgang mit Arzneimittel}



\section{Kurzbeschreibung wichtiger Notfallmedikamente}


\subsection{Herzrhythmus}

\opt{STATUS-EXPERIMENTAL}{
\Layout{57}{}{}{}{
\begin{tabularx}{\linewidth}{llccccX}
\hline

&

&
\multicolumn{3}{c|}{\bfseries
EKG} 
&
\bfseries Akz. Bdl. 
&
\bfseries Dosierung
\\\addlinespace\hline\addlinespace

&

&
\bfseries PR
&
\bfseries QRS
&
\bfseries QT 
&

&

\\\addlinespace\hline\addlinespace
\TabSub{Adenosin} 
&

&

&

&

&

&
Bolus: NW + KI: 

\\\addlinespace\hline\addlinespace
\bfseries Amiodaron III
&
~
&

&

&

&

&
Dos: 

\\\addlinespace\hline\addlinespace
\bfseries Digoxin 
&
~
&
${\uparrow}$
&
-
&
${\downarrow}$
&
beschl. 
&
Initial: 
0,5\Mg* über \VonBis{5}{30}\Min* \Abk{IV}, dann 0,25\Mg* alle
\VonBis{4}{6}\Hour* (max: \VonBis{1}{1,5}\Mg* am 1. Tag)

Erh:  
Je nach Kammerfrequenz und Nierenfunktion kann die \Abk{IV}-Gabe
wiederholt oder PO weitergefahren werden. 
\\\addlinespace\hline\addlinespace
\bfseries Flecainid 
&
~
&

&

&

&

&
Initial: 
 
Erh: 
\\\addlinespace\hline\addlinespace
\bfseries Lidocain Ib
&
~
&
-- 
&
-- 
&
-- 
&
-- 
&
\begin{compactdesc}
  \item[Dos] Pulslose KT oder Asystolie: 75-100~in 1-2~min \Abk{IV}(1-1,5\MgKg*),
  bei Bedarf alle 5~Min 50\Mg* in 1~Min~\Abk{iv} (0,5-0,75\MgKg*); max
  200-300\Mg* während der ersten Stunde
  \item[Hämodynamisch stabile KT:]  50\Mg* (0,5-0,75\MgKg*)~\Abk{iv} dann el
  CV, evt Erhaltung 1,5-5~g/24h~\Abk{iv}
\end{compactdesc}
\\\addlinespace\hline\addlinespace
\bfseries Procainamid Ia
&
~
&

&

&

&

&
Initial:  Erh:
\\\addlinespace\hline\addlinespace
\bfseries Propafenon Ic
&
~
&

&

&

&

&
Initial:  
\\\addlinespace\hline\addlinespace
\bfseries Propanolol II
&
~
&

&

&

&
~
&
Initial:  
\\\addlinespace\hline\addlinespace
\bfseries Sotalol III
&
~
&
${\uparrow}$
&
${\uparrow}$
&
${\uparrow}$
&
verl. 
&
 Initial:
 20\Mg* über 5~min IV; bei Bedarf nach 20~min wiederholen:
20\Mg* über 20~min; max 1,5\MgKg*

  Erh: \Abk{IV}: max 1,5~\MgKg* \Abk{IV}aufgeteilt in 4 Tagesdosen\newline
PO: 160-320\Mg* aufgeteilt in 2 Tagesdosen (max. 640\Mg/d)
\\\addlinespace\hline\addlinespace
\bfseries Verapamil IV
&

&
${\uparrow}$${\uparrow}$
&
--
&
--
&

&
 Initial:  
 5-10\Mg* über 10~min~\Abk{iv} (0,075-0,15\MgKg*), bei Bedarf nach
30~min wiederholen

Erh:
 \Abk{IV}: 0,005\MgKg*/min \Abk{IV}kont. \newline
PO: 3x~80-160\Mg/d (max 480\Mg/d)
\\\addlinespace\hline\addlinespace
\end{tabularx}
}
{Tabelle}
{}

} % STATUS-EXPERIMENTAL

\opt{STATUS-EXPERIMENTAL}{
* Adenosin, Verapamil und Propanolol (u.a. BB) beeinflussen die akzess
Bündel nicht. Bestehen aber solche Bündel (zB WPW-Syndrom),
können diese Substanzen maligne Arrhythmien auslösen (KT, KF),
indem die AV-Überleitung verlangsamt oder blockiert wird, was den
akzess Bündel die Möglichkeit gibt, den Ventrikel direkt zu
depolarisieren (präexzitieren).
} % STATUS-EXPERIMENTAL


\subsubsection{Tabelle Antiarrhythmika}

\opt{STATUS-EXPERIMENTAL}{
\Layout{57}{}{}{}{
\begin{tabulary}{\linewidth}{LLLL}
\hline
~
 &
1. Wahl
(stabil) &
2. Wahl &
Alternativen, Bemerkungen
\\\addlinespace\hline\addlinespace
 VHFlimmern &
~
 &
~
 &
~
\\\addlinespace\hline\addlinespace
\bfseries VHFlattern &
~
 &
~
 &
~
\\\addlinespace\hline\addlinespace
\EE{Reentry-Tachy:}

\EE{AV-Knoten Reentry Tachykardie}
(160-220/min)

\EE{ Anterograde Reentry Tachykardie bei akzess.
Bündel} (150-250/min)

\EE{schmale QRS!}
&
Vagale Manöver 
&
Adenosin 
&
Gute LVEF (nach abst. Indikation):


\begin{inparaenum}
  \item Verapamil oder Diltiazem
  \item BB
  \item Digoxin oder Procainamid
  \item Cardioversion
\end{inparaenum}

LVEF {\textless} 40\PC* oder chron HI

\begin{inparaenum}
  \item Amiodaron
  \item Digoxin
  \item Diltiazem (nicht Verapamil!)
  \item Kardioversion
\end{inparaenum}
\\\addlinespace\hline\addlinespace
\EE{Präexzitationssyndrome\textmd{ (zB WPW)}}

(= Retrograde Reentry Tachykardie bei akzess Bündel)

\EE{breites QRS!}
&
Procainamid IV

Flecainid \Abk{IV}
&
Amiodaron 
&
Keine Medikamente die den AV-Knoten
verlangsamen oder blockieren:

Digoxin, Adenosin, CA (Verapamil!),
BB
\\\addlinespace\hline\addlinespace
\EE{Vorhoftachykardie (uni-/multifokal) {\textpm} Block}

(VH: 100-250/min) 
&
~
 &
~
 &
~
\\\addlinespace\hline\addlinespace
Kammertachykardie 
&
~
 &
~
 &
~
\\\addlinespace\hline\addlinespace
~
 &
~
 &
~
 &
~
\\\addlinespace\hline
\end{tabulary}
}{Tabelle Antiarrhythmika}{}

Vorhof:

AV-Knoten:

Kammer:

} % STATUS-EXPERIMENTAL
%

%%%%%%%%%%%%%%%%%%%%%%%%%%%%%%%%%%%%%%%%%%%%%%%%%%%%%%%%%%%
%%%%%%%%%%%%%%%%%%%%%%%%%%%%%%%%%%%%%%%%%%%%%%%%%%%%%%%%%%%
\subsubsection{Ajmalin (I)}

\opt{STATUS-EXPERIMENTAL}{
Gilurytmal (10\Ml*/\,50\Mg, 2\Ml*/\,50\Mg*)

Wirkt auf VH u Ventrikel

\begin{labeling}[:]{mm}
  \item[Ind]
  \item[KI] Ventrikulare Tachykardiem Supraventrikuläre Tachykardie,
  WPW-Syndrom\footnote{Wolf-Parkinson-White-Syndrom},
  LGL-Syndrom\footnote{Lown-Gang-Levine-Syndrom}
  \item[KI]  Bradykardie, AV-Block, Schenkelblock, manifeste Herzinsuffizenz
  (ausgen. durch Arrhythmie)
  \item[Rx] Erw. 50\Mg* über mind. 5 Min \Abk{IV}

  Wh.: 50\Mg* nach 10 Min

  Perf.: \VonBis{0,5}{1}\MgKgH
  
  RR-Monitorkontrolle, Abbruch bei Verbreiterung der QRS-Komplexe,
  Temp.-labil
\end{labeling}
} % STATUS-EXPERIMENTAL
%%%%%%%%%%%%%%%%%%%%%%%%%%%%%%%%%%%%%%%%%%%%%%%%%%%%%%%%%%%
%%%%%%%%%%%%%%%%%%%%%%%%%%%%%%%%%%%%%%%%%%%%%%%%%%%%%%%%%%%
\subsubsection{Adenosin}
\opt{STATUS-EXPERIMENTAL}{

6\Mg\,/\,2\Ml

\begin{labeling}[:]{mm}
  \item[Ind] Diag. \& Therapie tachykarder Rhythmusstörungen:

  \EE{Reentrytachykardie} aus dem Bereich des des AV-Knotens wird
  terminiert

  \EE{SVT} wird verlangsamt

  \EE{VT}: kein Effekt
  \item[KI] \VonBis{2}{3}{\textdegree} AV-Block, Sick-Sinus-Syndrome

  \EE{CAVE}: bei Schwangerschaft, VH-Flattern/Flimmern, Asthma bronschiale
  \item[Rx] 6\Mg* rasch \Abk{IV}, bei 0 Erfolg Wiederholung 12\Mg, dann 18\Mg
  nach jeweils 60\Sek*.

  oder: \textbf{(3)-6-9-12} 
  
  \E{Kinder}: \VonBis{0,05}{0,15}\MgKg* in Schritten zu
  0,05\MgKg* in 2-Min-Abständen
  
  Dosisreduktion bei Z.\,n. Herz-TX (1/3-1/5) und bei \ce{Ca}-Antagonisten vom
  Nifedipin-Typ (1/4)

  Defi-Bereitschaft. Extrem kurze HWZ ({\textless}\,10 Sek). Es tritt
  ein \EE{AVB 3{\textdegree}} ein (kurz 0-Linie, keep cool)
\end{labeling}

\subparagraph{Antidot}: Theophyllin


} % STATUS-EXPERIMENTAL


%%%%%%%%%%%%%%%%%%%%%%%%%%%%%%%%%%%%%%%%%%%%%%%%%%%%%%%%%%%
%%%%%%%%%%%%%%%%%%%%%%%%%%%%%%%%%%%%%%%%%%%%%%%%%%%%%%%%%%%
\subsubsection{Amiodaron (\textrm{III})}
\opt{STATUS-EXPERIMENTAL}{

Sedacoron (1\,A\,/\,150\Mg*)
\begin{labeling}[:]{mm}
\item[KI]  bei Ventrikulärer Tachykardie keine KI, sonst:

Iodallergie, Sinusbrady, Hyperthyreose, K\,\Sign{↓}, Torsade
de pointes, Gravid., Stillen

\item[Ind] Reanimation bei KF/pulsl. VT:

\item[Rx]  Bolus 300\Mg* vor 4. Schock, überlege 1x Rep
150\Mg

\item[Ind] Therapie-resistente tachykarde supraventrikuläre oder ventrikuläre
Rhythmusstörung

\item[Rx] \MgKg* langsam über mind. 3 Min unter
Monitoring. Wiederholung frühestens nach 15\Min*.\newline
Einmalige Infusion: 300\Mg* in 250\Ml* G5\PC* innerh. v. \VonBis{0,5}{2} Std

\item[Ind] WPW

\item[Rx] \VonBis{0,4}{0,8}\MgKgH

ARREST- und ALIVE-Studie beweisen besseres klinisches Outcome vs. Placebo und
Lidocain

\Ca{Cave: \textbf{LQTS}, Bradykardie, Blöcke bis Asystolie, v.\,a. in Kombination
mit anderen Antiarrhythmika (\textbf{Ca-Blocker!!!}), Steigerung ICP}
\end{labeling}

} % STATUS-EXPERIMENTAL
%%%%%%%%%%%%%%%%%%%%%%%%%%%%%%%%%%%%%%%%%%%%%%%%%%%%%%%%%%%
%%%%%%%%%%%%%%%%%%%%%%%%%%%%%%%%%%%%%%%%%%%%%%%%%%%%%%%%%%%
\subsubsection{Atropin}

\Z{Atropinum sulfuricum (1\,A = 0,5\Mg*)}
\begin{labeling}[:]{mm}
  \item[Ind] Reanimation: PEA, Asystolie; Sinusbrady, Prophylaxe und Therapie.
  Vergiftung mit Alkylphosphaten der Hypersalivation
  \item[KI] Glaukom, Fieber, Hyperthyreose, Tachyarrhythmie,
  Prostatahypertrophie, Myasthenie
  \item[Rx] \begin{itemize}
    \item \EE{Reanimation}: Erw: 3\Mg* einmalig; \E{Kinder}:
    \item \EE{Hypersalivation}: \VonBis{0,01}{0,02}\MgKg* (od:
    0,5\Mg* Bolus, bis 3\Mg*)
    \item \EE{Alkylphosphat-Intox}: Initial \VonBis{2}{4}\Mg, dann 0,5\Mg* alle
    5\Min*.
  \end{itemize}
  \item[An] Endobronch. Anw. bei Reanimation möglich, paradoxer Effekt bei
  Unterdosierung, auch spasmolytischer und antiemetischer Effekt, Temp.-stabil
\end{labeling}

%%%%%%%%%%%%%%%%%%%%%%%%%%%%%%%%%%%%%%%%%%%%%%%%%%%%%%%%%%%
%%%%%%%%%%%%%%%%%%%%%%%%%%%%%%%%%%%%%%%%%%%%%%%%%%%%%%%%%%%
\subsubsection{Methyldigoxin}
\opt{STATUS-EXPERIMENTAL}{

\Z{Lanitop (1\Ml*/\,0,2\Mg*)}

\begin{labeling}[:]{mm}
  \item[Ind] VHflattern + Vhflimmern mit schneller
  Überleitung
  \item[Rx] Erw: \VonBis{0,1}{0,2}\Mg* (\VonBis{\nicefrac{1}{2}}{1} Amp zu
  0,2\Mg*)
  \item[KI]  Flykosidintox, AVB \VonBis{2}{3}{\textdegree}, Sinusbradykardie,
  Hypo-K\textsuperscript{+}, Hyper-Ca\textsuperscript{++}, vor Kardioversion
  \item[Cave] \Cave{Nicht bei bereits digitalisierten Patienten, nicht nach iv-Gabe von
  Calcium, nur wenn hämodynamisch wirksam.}
\end{labeling}

} % STATUS-EXPERIMENTAL

%%%%%%%%%%%%%%%%%%%%%%%%%%%%%%%%%%%%%%%%%%%%%%%%%%%%%%%%%%%
%%%%%%%%%%%%%%%%%%%%%%%%%%%%%%%%%%%%%%%%%%%%%%%%%%%%%%%%%%%
\subsubsection{Esmolol (\textrm{II})}
\opt{STATUS-EXPERIMENTAL}{

\Z{Brevibloc (10\Ml*/\,100\Mg, 10\Ml*/\,2500\Mg*)}

\begin{labeling}[:]{mm}
  \item[Ind]Tachyarrhythmie, insbes. SVT,
  Thyreotoxikose!, Flimmerprophylaxe nach MCI
  \item[KI]  Bradykardie, NYHA \VonBis{III}{IV}, (kardiogener) Schock, AVB
  \VonBis{2}{3}{\textdegree}, Sinusknotensyndrom, akuter Asthmaanfall
  \item[Rx] Bolus: 0,5\MgKg*, Rep. idem

  \E{Perfusor}: 3-18\MgKgH
  
  \EE{Monitor/RR-Kontrolle, langsam}
  
  CAVE: \textbf{keinesfalls mit Ca-Antagonisten (Verapamil) kombinieren}!
\end{labeling}
} % STATUS-EXPERIMENTAL
%%%%%%%%%%%%%%%%%%%%%%%%%%%%%%%%%%%%%%%%%%%%%%%%%%%%%%%%%%%
%%%%%%%%%%%%%%%%%%%%%%%%%%%%%%%%%%%%%%%%%%%%%%%%%%%%%%%%%%%
\subsubsection{Lidocain (2\PC*, \PC* x10=mg/ml) (\textrm{I})}
\begin{labeling}[:]{mm}
  \item[Ind] ischämiebed. ventr. Tachyarrh. (VT, VES),
  Breitkomplextachykardie beim MCI
  \item[Rx]Bolus 1\MgKg* langsam \Abk{IV}, dann Perf.
  \VonBis{2}{4}\MgKgH\newline
  tief endobronch. Bolus \VonBis{2}{3}\MgKg
  \item[KI] AV-Block \VonBis{2}{3}{\textdegree}, Bradykardie, dekomp.
  Herzinsuffizienz (ausgen. durch Arrhythmie)
  \Fazit{Bei Flimmern ist Amiodaron 1. Wahl. (auch bei Tachykardie)}
\end{labeling}



%%%%%%%%%%%%%%%%%%%%%%%%%%%%%%%%%%%%%%%%%%%%%%%%%%%%%%%%%%%
%%%%%%%%%%%%%%%%%%%%%%%%%%%%%%%%%%%%%%%%%%%%%%%%%%%%%%%%%%%
\subsubsection{Metoprolol}
\opt{STATUS-EXPERIMENTAL}{

} % STATUS-EXPERIMENTAL
%%%%%%%%%%%%%%%%%%%%%%%%%%%%%%%%%%%%%%%%%%%%%%%%%%%%%%%%%%%
%%%%%%%%%%%%%%%%%%%%%%%%%%%%%%%%%%%%%%%%%%%%%%%%%%%%%%%%%%%
\subsubsection{Orciprenalin}
\opt{STATUS-EXPERIMENTAL}{

\Z{Alupent (1\Ml*/\,0,5mg, 10ml)}
\begin{labeling}[:]{mm}
  \item[Ind] Bradykardie bei AVB 3{\textdegree}, Betablocker-Intox
  \item[Rx] Bolus 0,02\MgKg* (auch \Abk{SC}, \Abk{IM})

  Mit \ce{NaCl} verdünnt ml-weise unter Monitorkontrolle oder Perfusor geben.
  \item[KI] Tachykardie, Thyreotoxikose, hypertrophe obstruktive Kardiomyopathie
\end{labeling}
} % STATUS-EXPERIMENTAL
%%%%%%%%%%%%%%%%%%%%%%%%%%%%%%%%%%%%%%%%%%%%%%%%%%%%%%%%%%%
%%%%%%%%%%%%%%%%%%%%%%%%%%%%%%%%%%%%%%%%%%%%%%%%%%%%%%%%%%%
\subsubsection{Verapamil (\textrm{IV})}
\opt{STATUS-EXPERIMENTAL}{

\Z{Isoptin (2\Ml*/\,5\Mg, 20\Ml*/\,50\Mg*) }

Sinusknoten, AV-Überleitung

\begin{labeling}[:]{mm}
  \item[Ind]
  \item[KI] SVT,
  VHFlattern/Flimmern m schneller Überleitung
  \item[Rx] Erw: \VonBis{2,5}{5}\--\,10\Mg;
  Kind: 0,1\MgKg*;
  \E{Perfusor}: \VonBis{0,05}{0,1} \MgKgH (zB 50\Mg/50\Ml*/\,5)
  \item[KI]  AV-Block, Herzinsuffizenz, kardiogener Schock, Sick-Sinus-Syndrome,
  WPW-Syndrom, LGL-Syndrom, Kinder {\textless} 12\,J
  
  \Cave{Keinesfalls mit Beta-Blocker}
  
  Monitor-/RR-Kontrolle
\end{labeling}
} % STATUS-EXPERIMENTAL

%%%%%%%%%%%%%%%%%%%%%%%%%%%%%%%%%%%%%%%%%%%%%%%%%%%%%%%%%%%
%%%%%%%%%%%%%%%%%%%%%%%%%%%%%%%%%%%%%%%%%%%%%%%%%%%%%%%%%%%
\subsection{Katecholamine}

%%%%%%%%%%%%%%%%%%%%%%%%%%%%%%%%%%%%%%%%%%%%%%%%%%%%%%%%%%%
%%%%%%%%%%%%%%%%%%%%%%%%%%%%%%%%%%%%%%%%%%%%%%%%%%%%%%%%%%%
\subsubsection{Adrenalin}

\Layout{23}{Quick}
{
\begin{labeling}[:]{mm}
  \item[Präp] \Conc{1}{10\,000}: L-Adrenalin, \Conc{1}{1\,000} Suprarenin.
  \item[Ind] Reanimation (schockbarer und nicht-schockbarer Rhythmus),
  Kardiogener Schock, Anaphylaxie > 3{\textdegree},  schweres Asthma bronchiale.
  \item[KI]Im Notfall keine, ansonsten:
  Hypertonie, KHK, Glaukom, Tachyarrhythmie, Hyperthyreose,
  Phäochromocytom.
  \item[\textrecipe]
  \begin{itemize}
    \item \EE{Reanimation:} Erw: 1\Mg, \EE{\E{Kinder}: 0,01\MgKg}
    (1\Ml*/\,10\,kg (bei 1:10\,000)); Wiederholung alle 3\,min.)
    \item \EE{Kardiogener Schock:}
    \VonBis{0,025}{0,3}\,{\textmu}g\,/\,kgKG\,/\,min
    \item \EE{Anaphylaxie:} 3{\textdegree} und \EE{schweres Asthma bronchiale}: 
    \E{Erw.}: 0,1\Mg* \Abk{IV} oder \VonBis{0,2}{0,3}\Mg* \Abk{SC}/\Abk{IM}; 
    \E{Kinder}: 0,01\MgKg/\Abk{SC}/\Abk{IM}; 
    Wiederholung alle \VonBis{5}{15} Min
    \item \EE{Kruppanfall:} \Conc{1}{1\,000} (sic!) im Vernebler \Conc{1}{2} mit \ce{NaCl}
    verdünnt, 0,1\MgKg* inhal.
  \end{itemize}
  \item[Anm] Bei tiefer endobronchialer Applikation \VonBis{2}{3}-fache Dosis.
  Bei gleichz. Ca-Gabe \EE{Arrhythmien}; \EE{NaBic} über gl. Zugang inakt.
  Katech.; Kein Adrenalin in Akren spritzen; sehr Temp.-labil
  ({\textgreater}\,40\,{\textcelsius}), nur farblose Lösungen verwenden; Je höher die
  Dosis, desto schlechter neurologisches Outcome.
  \item[Cave] Bei Betablocker-Pat. überwiegend
  {$\alpha$}-Rezeptoren-Wirkung.
\end{labeling}
}
{}
{./icons/under-construction.png}
{}
{}









%%%%%%%%%%%%%%%%%%%%%%%%%%%%%%%%%%%%%%%%%%%%%%%%%%%%%%%%%%%
%%%%%%%%%%%%%%%%%%%%%%%%%%%%%%%%%%%%%%%%%%%%%%%%%%%%%%%%%%%
\subsubsection{Dobutamin}
\opt{STATUS-EXPERIMENTAL}{

\opt{STATUS-DRAFT}{{\FontExperimental
\begin{labeling}[:]{mm}
  \item[Beschreibung:]
  \item[Ind]
  \item[KI]
  \item[Präparate:]
  \item[Ind] Kardiogener Schock mit Rückwärtsversagen
  (Stauung, RR {\textgreater}\,90)
  \item[KI] Tachyarrhythmie, Dehydratation,
  Hypertonie, Sulfit-Allergie, (subvalvuläre Aortenstenose)
  \item[Rx] \VonBis{2}{10} (20) \MugKgMin .
  \textit{Infusion:} 250\Mg* in 500\Ml* RLH
  \VonBis{7}{18}\,--\,36\,\nicefrac{gtt}{\Min}
  
  Monitor/RR-Kontrolle, wenn möglich mit Spritzenpumpe,
  Temp.-labil
  
  \Ca{Cave: bei Beta-Blocker-Pat. überwiegend Alpha-Wirkung, Inaktivierung
  durch \textbf{NaBic}!}
\end{labeling}
}}
} % STATUS-EXPERIMENTAL
%%%%%%%%%%%%%%%%%%%%%%%%%%%%%%%%%%%%%%%%%%%%%%%%%%%%%%%%%%%
%%%%%%%%%%%%%%%%%%%%%%%%%%%%%%%%%%%%%%%%%%%%%%%%%%%%%%%%%%%
\subsubsection{Dopamin}
\opt{STATUS-EXPERIMENTAL}{

\Experimental{
\begin{labeling}[:]{mm}
  \item[Beschreibung:]
  \item[Ind]
  \item[KI]
  \item[Präparate:]

5\Ml*/\,50\Mg, 10\Ml*/\,200\Mg, 50\Ml*/\,250\Mg

\item[Ind] Kardiogener Schock mit Vorwärtsversagen
(RR {\textless}\,90), Nierenversagen, septischer Schock

\item[KI] Tachyarrhythmie, Dehydratation, Hypertonie,
Sulfit-Allergie, Hyperthyreose, Phäochromocytom

\item[Rx] \VonBis{2}{10} (20) {\textmu}g\,/\,kgKG\,/\,min

Monitor/RR-Kontrolle, wenn möglich mit Spritzenpumpe,
Temp.-labil
\end{labeling}
}
} % STATUS-EXPERIMENTAL
%%%%%%%%%%%%%%%%%%%%%%%%%%%%%%%%%%%%%%%%%%%%%%%%%%%%%%%%%%%
%%%%%%%%%%%%%%%%%%%%%%%%%%%%%%%%%%%%%%%%%%%%%%%%%%%%%%%%%%%

\subsubsection{Noradrenalin = Norepinephrin (NOR)}
\opt{STATUS-EXPERIMENTAL}{
\begin{labeling}[:]{mm}
  \item[Beschreibung:]
  \item[Ind]
  \item[KI]
  \item[Präparate:]
\end{labeling}


Arterenol (1\Ml*/\,1\Mg, 5\Ml, 25\Mg*)

\begin{labeling}[:]{mm}
  \item[Ind] Zustände mit niedrigem perivaskulärem Wiederstand
  (z.\,B. septischer Schock, Überdosierung von Vasodilatatoren,
  Therapie-resistente Hypotonie)

  \item[Rx] gro{\ss}e Variablilität

  \VonBis{0,01}{\EE{0,1}}\,--\,1,0\MugKgMin

  Erw: 5\Mg\,/\,50\Ml*/\,4,2\Ml*/\Hour*

  \E{Kinder}: pro 10\,kg: 5\Mg\,/\,50\Ml*/\,0,6\Ml*/\,	h
 \item[KI] Hypertonie, Tachykardie, Cave bei Inhalationsnarkotika,
 Cor pulmonale, KHK, Gravidität, Hyperthyreose, Glaukom

 invasive RR-Messung, Rekap kontrollieren
 \item[NW] Nekrosegefahr wenn peripher-venös od. para,
 Zentralisation, Lungenödem, Angst, Unruhe, Tachykardie, Bradykardie
\end{labeling}
%
} % STATUS-EXPERIMENTAL
\subsection{Kortikoide}

%

%%%%%%%%%%%%%%%%%%%%%%%%%%%%%%%%%%%%%%%%%%%%%%%%%%%%%%%%%%%
%%%%%%%%%%%%%%%%%%%%%%%%%%%%%%%%%%%%%%%%%%%%%%%%%%%%%%%%%%%
\subsubsection{Methyl-Prednisolon}
\opt{STATUS-EXPERIMENTAL}{

Solu-Medrol (TSA 250, 500, 1000\Mg*)
\begin{labeling}[:]{mm}
  \item[Beschreibung:]
  \item[Ind]
  \item[KI]
  \item[Präparate:]
\end{labeling}
\begin{labeling}[:]{mm}\item[Ind] \item[KI] Rückenmarkstrauma

\item[Rx] Bolus: 30\MgKg* in 15\,min, Erhaltung:
5,4\MgKgH über 23\Hour*(/47\Hour*) (NASCIS)

\item[KI]  I. Notfall SHT%
%ÖNK Konsensus 2005: 6. Die Behandlung mit hochdosiertem Methylprednisolon ist bei gleichzeitigem Vorliegen eines Schädel-Hirn
, sonst Überempfindlichkeit, Pilz-, Herpes-, TB-Infektion, ak.
Spychosen, Ulcera

Therapiebeginn innerhalb der ersten 8 Stunden nach Querschnittslähmung


Wirkung beim WS-Trauma umstritten (ÖNK Consensus
2005\footnote{http://www.agn.at/dokumente/solumedrol.pdf}:  präklinisch nicht
empf.), wenn, dann nur Methyl-Prednisolon!

NASCIS-Schema: Bolus (30\MgKg) nach {\textless}3\Hour*, spätestens
{\textless} 8\Hour*, gefolgt von Erhaltungs-Infusion 5,4\MgKgH für
23\Hour* (bzw. 47h, Rx erst ${\geq}$ 3\Hour* nach Trauma).


\end{labeling}
} % STATUS-EXPERIMENTAL
%%%%%%%%%%%%%%%%%%%%%%%%%%%%%%%%%%%%%%%%%%%%%%%%%%%%%%%%%%%
%%%%%%%%%%%%%%%%%%%%%%%%%%%%%%%%%%%%%%%%%%%%%%%%%%%%%%%%%%%
\subsubsection{Prednisolon}
\opt{STATUS-EXPERIMENTAL}{

Solu-Dacortin, Solu-Decortin (TSA: 25, 50, 250, 1000\Mg*)
\begin{labeling}[:]{mm}
  \item[Beschreibung:]
  \item[Ind]
  \item[KI]
  \item[Präparate:]


\textbf{\textit{\textcolor[rgb]{0.0,0.5019608,0.0}{In}}}\textbf{ \&
}\textbf{\textit{Rx:
}}\textcolor[rgb]{0.0,0.5019608,0.0}{Stoffwechselkoma:} \VonBis{50}{100}\Mg,
\textcolor[rgb]{0.0,0.5019608,0.0}{Status asthmaticus} \VonBis{250}{500}\Mg,
\textcolor[rgb]{0.0,0.5019608,0.0}{Anaphylaxie} \VonBis{1}{3}\Gramm*, (tox.
\textcolor[rgb]{0.0,0.5019608,0.0}{Lungenödem, Reizgasinhal.} \VonBis{1}{3}\Gramm*)

\item[KI]  Im Notfall keine, Überempfindlichket, Pilz- Herpes- Tb-Infektion,
ak Psychosen, Ulcera

Kortikoide haben keine Sofortwirkung, Notfallindikationen sind zT
umstritten
\end{labeling}
%
} % STATUS-EXPERIMENTAL

\subsection{Narkose}
\opt{STATUS-EXPERIMENTAL}{

\begin{labeling}[:]{mm}
  \item[Präp] ---.
  \item[Beschreibung] ---.
  \item[Ind] ---.
  \item[KI] ---.
  \item[Rx] ---.
  \item[Cave] ---.
  \item[Anm] ---.
\end{labeling}
} % STATUS-EXPERIMENTAL
Relaxanz


%%%%%%%%%%%%%%%%%%%%%%%%%%%%%%%%%%%%%%%%%%%%%%%%%%%%%%%%%%%
%%%%%%%%%%%%%%%%%%%%%%%%%%%%%%%%%%%%%%%%%%%%%%%%%%%%%%%%%%%
\subsubsection{Cis-Atracurium}
\opt{STATUS-EXPERIMENTAL}{

\begin{labeling}[:]{mm}
  \item[Präp] (2,5\Ml*/\,5mg, 5\Ml*/\,10mg, 10\Ml*/\,20\Mg*)
  \item[Beschreibung]
  \item[Ind] Relaxation, Präcurarisierung, TIVA
  \item[KI] Cave bei Myasthenia gravis, neuromusk. Erkr. und
Gravid.\newline
Langsame Appl. bei kardiovask. Risiken und allergischer Diathese
\item[Rx] Intub: 0,2\MgKg*; Rep: 0,03\MgKg\newline
\E{Perfusor}: \VonBis{3}{2}\,--\,1\MugKgMin
\item[Anm.:]Reines Isomer des Atracuriums, kaum Histaminausschüttung.
Wirk.-Eintritt \VonBis{120}{150}\,sek, Dauer \VonBis{40}{50}\,min

RepDosis wirkt etwa 20\,min.

Lagerung \VonBis{2}{8}{\textcelsius}
\end{labeling}

} % STATUS-EXPERIMENTAL

%%%%%%%%%%%%%%%%%%%%%%%%%%%%%%%%%%%%%%%%%%%%%%%%%%%%%%%%%%%
%%%%%%%%%%%%%%%%%%%%%%%%%%%%%%%%%%%%%%%%%%%%%%%%%%%%%%%%%%%
\subsubsection{Rocuronium}
\opt{STATUS-EXPERIMENTAL}{


\begin{labeling}[:]{mm}
  \item[Präp] \Z{\T{Esmeron}{\texttrademark} (50\Mg\,/\,5\Ml,
  100\Mg\,/\,10\Ml)}
  \item[Beschreibung] ---.
  \item[Ind] Relaxation, Präcurarisierung, TIVA.
  \item[KI] Myasthenia gravis, Lambert Eaton Sy., Zn.
  Poliomyelitis, strenge Indikationsstellung bei Grav./Stillen.
  \item[Rx] Intubation: \VonBis{0,3}{0,6}\MgKg*; Rep.:
  0,15\MgKg\newline Perf.: 0,3-0,6\MgKg\,/\Hour*.
  \item[Anm.] Bei \textit{0,6\MgKg} gute Intubationsbedingungen schon nach
  1\,min, Wirkdauer ca. \VonBis{40}{50}{\textquotesingle}. Wiederholung wirkt
  etwa \VonBis{13}{20}\,min.
\end{labeling}

} % STATUS-EXPERIMENTAL

%%%%%%%%%%%%%%%%%%%%%%%%%%%%%%%%%%%%%%%%%%%%%%%%%%%%%%%%%%%
%%%%%%%%%%%%%%%%%%%%%%%%%%%%%%%%%%%%%%%%%%%%%%%%%%%%%%%%%%%
\subsubsection{Succinylcholin}
\opt{STATUS-EXPERIMENTAL}{
\begin{labeling}[:]{mm}
  \item[Präp] Lysthenon (5\Ml*/\,100mg (2\PC*), TSA: 500\Mg* + 25\Ml* NaCl~0,9\PC* = 2\PC*)
  \item[Beschreibung]
  \item[Ind] Crash-Intubation
  \item[KI] Cholinesterasemangel, MH, Hyper-K+, neuromusk
Erkr, LuÖdem, perf Augenverl, schwerer LPS
\item[Anm.] \item[Rx] \VonBis{1}{1,5}~\MgKg*;
Präcurarisieren!


Wirkungseintr \VonBis{45}{60}\,sek, Wirkdauer \VonBis{3}{4}\,min.
\item[Cave] Cave bei Kindern und Jugendlichen; irrev Herzstillstände in Pat mit
neuromusk Erkr
\end{labeling}

} % STATUS-EXPERIMENTAL




%%%%%%%%%%%%%%%%%%%%%%%%%%%%%%%%%%%%%%%%%%%%%%%%%%%%%%%%%%%
%%%%%%%%%%%%%%%%%%%%%%%%%%%%%%%%%%%%%%%%%%%%%%%%%%%%%%%%%%%
\subsubsection{Diazepam}
\opt{STATUS-EXPERIMENTAL}{
\begin{labeling}[:]{mm}
  \item[Präp] Amp: Gewacalm (2\Ml*/\,10\Mg*), Gtt: Psychopax Drops,
  Rectal: 5\Mg, 10\Mg
  \item[Beschreibung]
  \item[Ind] Sedierung, Epilepsie, Eklampsie, Prämediaktion (CV), Rectal bei
  Krämpfen
  \item[KI] chron. Hyperkapnie, Myasthenia gravis, Intox mit
zentral dämpfenden Substanzen
\item[Rx i.V.]  0,2-0,5\MgKg*, maximal 20\Mg
\item[Cave] CAVE: Kumulation. Nicht bei Kindern {\textless}2\,a, auch
imAtemdepression
\item[Rx rect.] $\leq$ 15\,kg: 5\Mg, {\textgreater}15\,kg 10\Mg
\item[Anm.] Antidot: Flumazenil (Anexate)

\end{labeling}
} % STATUS-EXPERIMENTAL


%%%%%%%%%%%%%%%%%%%%%%%%%%%%%%%%%%%%%%%%%%%%%%%%%%%%%%%%%%%
%%%%%%%%%%%%%%%%%%%%%%%%%%%%%%%%%%%%%%%%%%%%%%%%%%%%%%%%%%%
\subsubsection{Flumazenil}
\opt{STATUS-EXPERIMENTAL}{

\begin{labeling}[:]{mm}
  \item[Präp] Anexate (5\Ml*/\,0,5mg, 10\Ml*/\,1\Mg*)
  \item[Beschreibung] Benzoantagonist
  \item[Ind] \Sign{↑} ICP, Epi, Intox mit Trizyklika
  \item[KI]
  \item[Rx.] initial \VonBis{0,2}{0,3}\Mg, Wiederholung 0,2\Mg* alle 60\,sek,
  Gesamtdosis: \VonBis{1}{2}\Mg. \E{Perfusor}: \VonBis{0,1}{0,4}\Mg/h;
  \E{Kinder}: \VonBis{4}{20}\,\nicefrac{{\textmu}g}{kg}
  \item[Anm.] evt. Auslösung von Entzugssymptomen, Temp.-stabil
\end{labeling}

} % STATUS-EXPERIMENTAL

\opt{STATUS-EXPERIMENTAL}{
\Layout{22}{Lorazepam}
{
\begin{labeling}[:]{mm}
  \item[Präp] Tavor, Temesta (1\Ml*/\,4\Mg*)
  \item[Beschreibung:]
  \item[Ind] Status epilepticus (1.W), Sedierung
  \item[KI] chronische Hyperkapnie, Myasthenia gravis, Intox. mit
  zentralal dämpfenden Substanzen
  \item[Rx.] Erw: \VonBis{2}{8}\Mg; Status: 4\Mg, Wiederholung mit 4\Mg* nach
  10\,min\newline \E{Kinder}: 0,1\MgKg
  \item[Anm.] Atemdepression, Antagonist: Anexate, Bei Erfolglosigkeit im Status
  Valproat
\end{labeling}
}
{}
{./icons/under-construction.png}
{}
{}

} % STATUS-EXPERIMENTAL




%%%%%%%%%%%%%%%%%%%%%%%%%%%%%%%%%%%%%%%%%%%%%%%%%%%%%%%%%%%
%%%%%%%%%%%%%%%%%%%%%%%%%%%%%%%%%%%%%%%%%%%%%%%%%%%%%%%%%%%
\opt{STATUS-EXPERIMENTAL}{
\Layout{22}{Midazolam}
{
\begin{labeling}[:]{mm}
  \item[Präp] Dormicum (1\Ml*/\,5mg, 3\Ml*/\,15mg, 5\Ml*/\,5mg (!!), 10\Ml*/\,50\Mg*)
  \item[Beschreibung]
  \item[Ind] Sedierung, zerebraler Krampfanfall,
  Status epi, Narkoseeinleitungchron, Hyperkapnie, Myasthenia gravis,
  ImzdS
  \item[KI]
  \item[Rx] Sedierung: \VonBis{0,015}{0,03} \MgKg\newline
  Krampf: \VonBis{0,05}{0,1}\MgKg\newline
  Status: bis 0,2\MgKg\newline
  Einleitung: \VonBis{0,1}{0,2}\MgKg
  \item[Anm.] Atemdepression gut steuerbar, Wirkdauder individuell
  \item[Cave]  paradoxe Wirkung va bei älteren
\end{labeling}
}
{}
{./icons/under-construction.png}
{}
{}
} % STATUS-EXPERIMENTAL

\opt{STATUS-EXPERIMENTAL}{
\Layout{22}{}
{

}
{}
{./icons/under-construction.png}
{}
{}
} % STATUS-EXPERIMENTAL
%%%%%%%%%%%%%%%%%%%%%%%%%%%%%%%%%%%%%%%%%%%%%%%%%%%%%%%%%%%
%%%%%%%%%%%%%%%%%%%%%%%%%%%%%%%%%%%%%%%%%%%%%%%%%%%%%%%%%%%
\subsubsection{Opiate}
\opt{STATUS-EXPERIMENTAL}{

} % STATUS-EXPERIMENTAL
%%%%%%%%%%%%%%%%%%%%%%%%%%%%%%%%%%%%%%%%%%%%%%%%%%%%%%%%%%%
%%%%%%%%%%%%%%%%%%%%%%%%%%%%%%%%%%%%%%%%%%%%%%%%%%%%%%%%%%%
\subsubsection{Fentanyl}
\opt{STATUS-EXPERIMENTAL}{




\begin{labeling}[:]{mm}
  \item[Präp] (2\Ml*/\,0,1mg=100{\textmu}g, 10\Ml*/\,0,5\Mg*)
  \item[Beschreibung:]
  \item[Ind] Opioidanalg., Einl. und Vert. d. Narkose
  \item[KI] Im Notfall keine, sonst: Stillen, Sgl., Asthma, akuter hepatischer
  Porphyrie, Myasthenia gravis, Geburtshilfe (Sektio)
  \item[Rx] Erw: \VonBis{0,05}{0,2}\Mg; \E{Kinder}: \VonBis{1}{2}\,(--\,5)
  {\textmu}g/kg \newline Dosisred. bei LPS, NINS, aälteren, schlehctem AZ.
  \item[Anm.] Tip: 0,1\,\Ml*/\,5\,{\textmu}g. {\textmu}-Agonist, Atemdepress.,
  insbes. bei Komb. mit Benzos. In höheren Dosen Thoraxrigidität. Antagonist:
  Naloxon.
\end{labeling}

} % STATUS-EXPERIMENTAL



%%%%%%%%%%%%%%%%%%%%%%%%%%%%%%%%%%%%%%%%%%%%%%%%%%%%%%%%%%%
%%%%%%%%%%%%%%%%%%%%%%%%%%%%%%%%%%%%%%%%%%%%%%%%%%%%%%%%%%%
\subsubsection{Morphin (hydrochlorid)}
\opt{STATUS-EXPERIMENTAL}{

Vendal (1\Ml*/\,10\Mg*)

\begin{labeling}[:]{mm}
  \item[Präp] ---.
  \item[Beschreibung] ---.
  \item[Ind] Analgesie insbes. Herzinfarkt.
  \item[KI]  im Notfall: Asthma, spastische Schmerzen, Sgl, Stillen.
  \item[Rx] Erw: \VonBis{2,5}{10}\Mg* \Abk{IV}, \E{Kinder}: 0,05\MgKg* iv; auch
  sc/im.
  \item[Cave] ---.
  \item[Anm] {\textmu}-Agonist, Atemdepression, parasympmimetische NW.
\end{labeling}

} % STATUS-EXPERIMENTAL





%%%%%%%%%%%%%%%%%%%%%%%%%%%%%%%%%%%%%%%%%%%%%%%%%%%%%%%%%%%
%%%%%%%%%%%%%%%%%%%%%%%%%%%%%%%%%%%%%%%%%%%%%%%%%%%%%%%%%%%
\subsubsection{Naloxon}
\opt{STATUS-EXPERIMENTAL}{



\begin{labeling}[:]{mm}
  \item[Präp]  Narcanti (1\Ml*/\,0,4\Mg*).
  \item[Beschreibung] Opiat-Antagonist .
  \item[Ind] ---.
  \item[KI] Überempfiundlichkeit, Atemdepression die nicht durch Opiate
ausgelöst wurde.
  \item[Rx] Erw: 0,4\Mg* \Abk{IV}, \Abk{IM}, sc\newline
\E{Kinder}: 0,1\MgKg* \Abk{IV}, \Abk{IM}, sc.
  \item[Cave]  Cave bei Pat mit vorbestehender Herzerkr. .
  \item[Anm]  {\textmu}-Antagonist, daher Entzugssympt., auch bei Gebärenden
indiziert---.
\end{labeling}

} % STATUS-EXPERIMENTAL



%%%%%%%%%%%%%%%%%%%%%%%%%%%%%%%%%%%%%%%%%%%%%%%%%%%%%%%%%%%
%%%%%%%%%%%%%%%%%%%%%%%%%%%%%%%%%%%%%%%%%%%%%%%%%%%%%%%%%%%
\subsubsection{Piritramid}
\opt{STATUS-EXPERIMENTAL}{

\begin{labeling}[:]{mm}
  \item[Präp]  Dipidolor (2\Ml*/\,15\Mg*).
  \item[Beschreibung] Opiatanalgetikum mit langer Wirkdauer.
  \item[Ind] ---.
  \item[KI]  Sgl, ak hepat Porphyrie, MAO-Hemmer .
  \item[Rx]  Erw: \VonBis{1}{2}\Ml* \Abk{IV}, \VonBis{2}{4}\Ml* im\newline
KK: 0,1\MgKg*, SK:
\A{???}\VonBis{0,1}{0,2}\MgKg\newline Wiederholung nach 30 Min .
  \item[Cave] ---.
  \item[Anm]  {\textmu}-Agonist, Wirkdauer ca 6\Hour*, Atemdepression, Dosisred.
  Bei älteren Pat, LPS, schlechtem AZ .
\end{labeling}

} % STATUS-EXPERIMENTAL

\subsection{Sonstige}

%%%%%%%%%%%%%%%%%%%%%%%%%%%%%%%%%%%%%%%%%%%%%%%%%%%%%%%%%%%
%%%%%%%%%%%%%%%%%%%%%%%%%%%%%%%%%%%%%%%%%%%%%%%%%%%%%%%%%%%
\subsubsection{Ketamin, S-}
\opt{STATUS-EXPERIMENTAL}{

\begin{labeling}[:]{mm}
  \item[Präp]  Ketanest-S (50\Mg\,/\,2\Ml, 25\Mg\,/\,5\Ml,
  250\Mg\,/\,10\Ml
  \item[Beschreibung] ---.
  \item[Ind] Analgesie, Anästhesie, Therapierefr. Bronchospasmus.
  \item[KI] erhöhter ICP\footnote{ICP: intrakranieller Druck} (relativ),
  penetrierende Augenverletzng, kard. Dekompensation, KHK.
  \item[\textrecipe :]
  Analgesie: \VonBis{0,1}{0,25}\MgKg* \Abk{IV},
  \VonBis{0,25}{0,5}\MgKg* \Abk{IM}; 
  Anästhesie: \VonBis{1}{2,5}\MgKg* \Abk{IV}, 2,5\,-\,4\MgKg* \Abk{IM};
  Bronchospasmus: \VonBis{1}{2,5}\MgKg* \Abk{IV}.
  \item[Cave] CAVE: Unterschiedliche Konzentrationen im Handel! .
  \item[Anm] Schluckreflexe / Spontanatmung \Sign{↑}, Rep: halbe Initialdosis.
  V.a. bei Kindern mit Atropin kombinieren (Hypersalivation), Komb. mit Benzos
  empfohlen.
\end{labeling}

} % STATUS-EXPERIMENTAL

%%%%%%%%%%%%%%%%%%%%%%%%%%%%%%%%%%%%%%%%%%%%%%%%%%%%%%%%%%%
%%%%%%%%%%%%%%%%%%%%%%%%%%%%%%%%%%%%%%%%%%%%%%%%%%%%%%%%%%%
\begin{labeling}[:]{mm}
  \item[Präp] ---.
  \item[Beschreibung] ---.
  \item[Ind] ---.
  \item[KI] ---.
  \item[Rx] ---.
  \item[Cave] ---.
  \item[Anm] ---.
\end{labeling}
\subsubsection{Etomidat}
\opt{STATUS-EXPERIMENTAL}{

\begin{labeling}[:]{mm}
  \item[Präp]  Hypnomidate (10\Ml*/\,20\Mg*).
  \item[Beschreibung] Kurzhypnotikum zur Einleitung und Intub..
  \item[Ind] ---.
  \item[KI]  Im Notfall keine, sonst: Neugeborene und Sgl bis 6 Monate, strenge
Ind. ind Gravid .
  \item[Rx]  \VonBis{10}{20}\,--\,40\Mg* (\VonBis{1}{2}\,A), max 60\Mg, \E{Kinder}:
  \VonBis{0,15}{0,30}\MgKg* .
  \item[Cave] ---.
  \item[Anm]  Bei rascher Injektion Masseterkrämpfe. Stillpause: 24\Hour* .
\end{labeling}

} % STATUS-EXPERIMENTAL

\subsubsection{Nalbuphin}
\opt{STATUS-EXPERIMENTAL}{

\begin{labeling}[:]{mm}
  \item[Präp] \Z{Nubain (1\Ml*/\,10mg, 2\Ml*/\,20\Mg*)} .
  \item[Beschreibung]  Analgetikum f. mittelstarke Schmerzen .
  \item[Ind] ---.
  \item[KI]  Atemstörungen, \Sign{↑} ICP, Grav, Sulfitallergie .
  \item[Rx]  \VonBis{0,1}{0,2}\MgKg*, max. 0,25\MgKg
  (Ceiling-Effekt) .
  \item[Cave] ---.
  \item[Anm]  {\textmu}-Agonist-Antagonist, daher Auslösung von Entzungssyndromen
bei süchtigen und deren Neugeborenen. .
\end{labeling}

} % STATUS-EXPERIMENTAL

\subsubsection{Propofol 1\PC*, 2\PC*}
\opt{STATUS-EXPERIMENTAL}{

\begin{labeling}[:]{mm}
  \item[Präp] \Z{(20\Ml*/\,200mg, 50\Ml*/\,500mg, 50\Ml*/\,1000mg !)} .
  \item[Beschreibung] ---.
  \item[Ind] Einleitung.
  \item[KI] relativ: Hypovolämie, Grav, Stillen, Geburtshilfe,
  Fettstoffwechselstörung, {\textless}1 Mon (1\PC*), {\textless}3a (2\PC*).
  \item[Rx]
  \begin{inparaitem}
  \item Einleitung: \VonBis{2,5}{5}\MgKg*, Erw:
  \VonBis{2}{3}\MgKg*, Ältere; \VonBis{1}{1,5}\MgKg*,
  (1\PC*: \VonBis{1}{5}\Ml*/\.10\,kg).
  \item TIVA: Beginn: 12{\MgKgH}, Reduz:
  \VonBis{6}{10}{\MgKgH} (1\PC*:
  \VonBis{6}{12}\Ml*/\,10\,kg\,/\Hour*)
  \item Sedierung: \VonBis{1}{4}{\MgKgH} (1\PC*:
  \VonBis{1}{4}\Ml/10kg/h).
  \end{inparaitem}
  \item[Cave] Hypotonie.
  \item[Anm]  Kardiodepressiv! ICP-Senkung.
\end{labeling}
} % STATUS-EXPERIMENTAL

%%%%%%%%%%%%%%%%%%%%%%%%%%%%%%%%%%%%%%%%%%%%%%%%%%%%%%%%%%%
%%%%%%%%%%%%%%%%%%%%%%%%%%%%%%%%%%%%%%%%%%%%%%%%%%%%%%%%%%%
\subsubsection{Thiopental }
\opt{STATUS-EXPERIMENTAL}{

\begin{labeling}[:]{mm}
  \item[Präp] ---.
  \item[Beschreibung] ---.
  \item[Ind] ---.
  \item[KI] ---.
  \item[Rx] ---.
  \item[Cave] ---.
  \item[Anm] ---.
\end{labeling}



{
Neostigmin nicht ZNS-gängig, Phyostigmin schon}


} % STATUS-EXPERIMENTAL


\subsection{Atmung}

%%%%%%%%%%%%%%%%%%%%%%%%%%%%%%%%%%%%%%%%%%%%%%%%%%%%%%%%%%%
%%%%%%%%%%%%%%%%%%%%%%%%%%%%%%%%%%%%%%%%%%%%%%%%%%%%%%%%%%%
\subsubsection{Fenoterol}
\opt{STATUS-EXPERIMENTAL}{

\begin{labeling}[:]{mm}
  \item[Präp] \Z{
  Berotec Dae (0,2mg/Hub); Partusisten (1\Ml*/\,0,025mg, 1\Ml*/\,0,5\Mg*)}.
  \item[Beschreibung] ---.
  \item[Ind] ---.
  \item[KI]  Dae: im Notfall keine. Tokolyse \Abk{IV}: starke Blutung, vorz
  Plazentalösung, Tachyarrhythmie, Thyreotoxikose, hypertrophe obstruktive
  CMP---.
  \item[Rx]
  \begin{inparaitem}
  \item Asthma Bronchiale: \VonBis{1}{2} Hübe, Wiederholung nach
  \VonBis{5}{10}{\textquotesingle}
  \item Tokolyse mit Dosieraerosol: \VonBis{2}{5} Hübe
  \item Tokolyse \Abk{IV}: sehr langsamer Bolus: 25\,{\textmu}g.
  \E{Perfusor}: \VonBis{0,01}{0,05}\,{\textmu}g\,/\,kgKG\,/\,min., zB
  2\Mg\,/\,50\Ml*/\,1-5 
  \end{inparaitem}.
  \item[Cave] ---.
  \item[Anm] DAe: Im Anfall meist unwirksam. \Abk{IV}: Kontrolle von Herzfrequenz von Mutter u
  Kind. Optional Hexoprenalin (Gynipral)..
\end{labeling}

} % STATUS-EXPERIMENTAL

%%%%%%%%%%%%%%%%%%%%%%%%%%%%%%%%%%%%%%%%%%%%%%%%%%%%%%%%%%%
%%%%%%%%%%%%%%%%%%%%%%%%%%%%%%%%%%%%%%%%%%%%%%%%%%%%%%%%%%%
\subsubsection{Terbutalin}
\opt{STATUS-EXPERIMENTAL}{

\begin{labeling}[:]{mm}
  \item[Präp]  \Z{Bricanyl 1\Ml*/\,0,5mg}.
  \item[Beschreibung] ---.
  \item[Ind] ---.
  \item[KI] ---.
  \item[Rx] ---.
  \item[Cave] ---.
  \item[Anm] ---.
\end{labeling}



\begin{labeling}[:]{mm}\item[Ind] \item[KI] Asthmaanfall

\item[Rx] 0,005\MgKg* (1A/100kg) \Abk{IV}/\Abk{SC}

\begin{labeling}[:]{mm}\item[Ind] \item[KI]\end{labeling} Tokolyse

\item[Rx] 0,25\Mg* ({\textonehalf} Amp)

\item[KI]  Cave bei Tachy, Thyreotox, hypertr obstr CMP, KHK, MCI, (subvalv)
Aortenstenose

Bei \Abk{IV}-Appl verdünnen, besser DT oder Perf. \Abk{SC} unverdünnt.
Monitor-/RR-Kontrolle.
\end{labeling}
} % STATUS-EXPERIMENTAL

%%%%%%%%%%%%%%%%%%%%%%%%%%%%%%%%%%%%%%%%%%%%%%%%%%%%%%%%%%%
%%%%%%%%%%%%%%%%%%%%%%%%%%%%%%%%%%%%%%%%%%%%%%%%%%%%%%%%%%%
\subsubsection{Theophyllin}
\opt{STATUS-EXPERIMENTAL}{

\begin{labeling}[:]{mm}
  \item[Präp] ---.
  \item[Beschreibung] ---.
  \item[Ind] ---.
  \item[KI] ---.
  \item[Rx] ---.
  \item[Cave] ---.
  \item[Anm] ---.
\end{labeling}

} % STATUS-EXPERIMENTAL

\subsection{Antipsychotika}

%%%%%%%%%%%%%%%%%%%%%%%%%%%%%%%%%%%%%%%%%%%%%%%%%%%%%%%%%%%
%%%%%%%%%%%%%%%%%%%%%%%%%%%%%%%%%%%%%%%%%%%%%%%%%%%%%%%%%%%
\subsubsection{Haloperidol}
\opt{STATUS-EXPERIMENTAL}{

\begin{labeling}[:]{mm}
  \item[Präp]  \Z{Halodol (1\Ml*/\,5\Mg*)}.
  \item[Beschreibung] ---.
  \item[Ind] Psychosen, Alk-intox, Delirium tremens.
  \item[KI] Depression zebtralnervöser Funktionen, Kinder, Parkinson,
malignes Neuroleptika-Sy.
  \item[Rx]  \VonBis{5}{10}\Mg* \Abk{IV}/\Abk{IM}.
  \item[Cave] Long-QT-Syndrom.
  \item[Anm] Dyskinesien, Temp.-stabil.
\end{labeling}

 Antiemetikum
\begin{labeling}[:]{mm}
\item[Rx] \VonBis{0,5}{1,0}\Mg* \Abk{IV}

\Cave{CAVE: }
\end{labeling}

} % STATUS-EXPERIMENTAL

%%%%%%%%%%%%%%%%%%%%%%%%%%%%%%%%%%%%%%%%%%%%%%%%%%%%%%%%%%%
%%%%%%%%%%%%%%%%%%%%%%%%%%%%%%%%%%%%%%%%%%%%%%%%%%%%%%%%%%%
\subsubsection{Psyquil (TM)}
\opt{STATUS-EXPERIMENTAL}{

\begin{labeling}[:]{mm}
  \item[Präp] ---.
  \item[Beschreibung] ---.
  \item[Ind] ---.
  \item[KI] ---.
  \item[Rx] ---.
  \item[Cave] ---.
  \item[Anm] ---.
\end{labeling}

} % STATUS-EXPERIMENTAL

%%%%%%%%%%%%%%%%%%%%%%%%%%%%%%%%%%%%%%%%%%%%%%%%%%%%%%%%%%%
%%%%%%%%%%%%%%%%%%%%%%%%%%%%%%%%%%%%%%%%%%%%%%%%%%%%%%%%%%%
\subsection{NSAR und Anti-''Au!''}

\subsubsection{Metamizol}
\opt{STATUS-EXPERIMENTAL}{
\begin{labeling}[:]{mm}
  \item[Präp] \Z{Novalgin (2\Ml*/\,1000\Mg, 5\Ml*/\,1500\Mg*)} .
  \item[Beschreibung] ---.
  \item[Ind]  mittlere bis schwere Kolikartige Schmerzen, Paracetamol-resistentes
Fieber .
  \item[KI] Porphyrie, Glucose-6-Phosphat-Dehydrogenasemangel, Grav.
  \item[Rx] Erw: \VonBis{0,5}{1}\Gramm* \Abk{IV}langsam!!\newline
Ki; 15\MgKg*\Abk{IV}langsam !!\newline
Rep nach 4 Std.
  \item[Cave] ---.
  \item[Anm] NW: RR-Abfall, roter Urin, verstärkte Blutungsneigung, Anaphylaxie,
Agranulacytose !!.
\end{labeling}

} % STATUS-EXPERIMENTAL

%%%%%%%%%%%%%%%%%%%%%%%%%%%%%%%%%%%%%%%%%%%%%%%%%%%%%%%%%%%
%%%%%%%%%%%%%%%%%%%%%%%%%%%%%%%%%%%%%%%%%%%%%%%%%%%%%%%%%%%
\subsubsection{Methylergometrin}
\opt{STATUS-EXPERIMENTAL}{

\begin{labeling}[:]{mm}
  \item[Präp] ---.
  \item[Beschreibung] ---.
  \item[Ind] ---.
  \item[KI] ---.
  \item[Rx] ---.
  \item[Cave] ---.
  \item[Anm] ---.
\end{labeling}

} % STATUS-EXPERIMENTAL

%%%%%%%%%%%%%%%%%%%%%%%%%%%%%%%%%%%%%%%%%%%%%%%%%%%%%%%%%%%
%%%%%%%%%%%%%%%%%%%%%%%%%%%%%%%%%%%%%%%%%%%%%%%%%%%%%%%%%%%
\subsubsection{Paracetamol}
\opt{STATUS-EXPERIMENTAL}{

\begin{labeling}[:]{mm}
  \item[Präp] ---
  \item[Beschreibung] ---
  \item[Ind]  Fieber, leichte bis mittelstarke Schmerzen .
  \item[KI] Überempf, G-6-P-DHG-Mangel, bei schwerem LPS und NINS Dosis
reduzieren, vor der 44. SSW (unreife Leber).
  \item[Rx] 15\MgKg* (Sgl: 125, KK: 250, Sk: 500, Erw: 1000\Mg*
  Supp).
  \item[Cave] ---
  \item[Anm] Antidot: N-Acetylcystein: 140\MgKg*, dann
  70\MgKg* alle 8\Hour*.
\end{labeling}

} % STATUS-EXPERIMENTAL

\subsection{Sonstige}


%%%%%%%%%%%%%%%%%%%%%%%%%%%%%%%%%%%%%%%%%%%%%%%%%%%%%%%%%%%
%%%%%%%%%%%%%%%%%%%%%%%%%%%%%%%%%%%%%%%%%%%%%%%%%%%%%%%%%%%
\subsubsection{4-DMAP (Dimethylaminophenol)}
\opt{STATUS-EXPERIMENTAL}{


\begin{labeling}[:]{mm}
  \item[Präp] ---
  \item[Beschreibung] ---
  \item[Ind]  Intox.: Cyanide, Blausäure,
Schwefelwasserstoff.
  \item[KI] Sulfit-Unverträglichkeit (allerg. Asthma), Mischintox mit \ce{CO}
  \item[Rx] initial \VonBis{3}{4}\MgKg*, anschließend
  \VonBis{100}{500}\Mg* Na-thiosulfat.
  \item[Cave] ---.
  \item[Anm] bei Überdosierung (Methämoglobinbildner): Toludinblau: 2\MgKg.
\end{labeling}
} % STATUS-EXPERIMENTAL
\section{Häufige Vor- und Dauermedikationen von Patienten}

\opt{STATUS-EXPERIMENTAL}{} % STATUS-EXPERIMENTAL

\ChapterEnd
